% Generated draft: requires mathematical and source review.
\documentclass[11pt]{article}
\usepackage[margin=25mm]{geometry}
\usepackage{fontspec}
\setmainfont{Latin Modern Roman}
\usepackage{amsmath,amssymb,mathtools,unicode-math}
\setmathfont{Latin Modern Math}
\usepackage{enumitem,needspace,xurl,hyperref}
\setlength{\parindent}{0pt}
\setlength{\parskip}{5pt}
\setlength{\emergencystretch}{3em}
\newcommand{\N}{\mathbb{N}}
\newcommand{\Z}{\mathbb{Z}}
\newcommand{\Q}{\mathbb{Q}}
\newcommand{\R}{\mathbb{R}}
\newcommand{\C}{\mathbb{C}}
\newcommand{\F}{\mathbb{F}}
\newcommand{\A}{\mathbb{A}}
\newcommand{\PP}{\mathbb{P}}
\newcommand{\E}{\mathbb{E}}
\newcommand{\eps}{\varepsilon}
\newcommand{\dd}{\,\mathrm{d}}
\begin{document}
\section*{UnsolvedMath ID: 30004079}
\textbf{Source problem code:} catalogue entry
\subsection{Definitions and mathematical statement}
Let \(R\) be a commutative Noetherian domain with identity and characteristic \(p>0\). The ring is \(F\)-finite if its Frobenius map \(F:R\to R\), \(r\mapsto r^p\), is finite; equivalently, \(R\), with scalar multiplication \(r\cdot x=r^p x\), is a finitely generated \(R\)-module. Let \(0\ne\mathfrak a\subseteq R\) be an ideal, allowing \(\mathfrak a=R\).

For a precise formulation, use the customary interpretation \(t\in\R_{>0}\). The supplied record does not explicitly specify the number domain of \(t\) or define its test-ideal conventions; the following definitions state the interpretation used in this draft.

For \(e\geq0\), put \(q=p^e\). Let \(R^{(e)}\) be the \(R\)-bimodule whose underlying additive group is \(R\), with ordinary left multiplication and right action \(x\cdot r=xr^q\). For an \(R\)-module \(M\), set \(F_R^e(M)=R^{(e)}\otimes_R M\). Define the \(\mathfrak a^t\)-tight closure of zero by
\[
0_M^{*\mathfrak a^t}
=
\left\{z\in M:\begin{array}{l}
\text{there exist }c\in R\setminus\{0\}\text{ and }e_0\geq0\text{ such that}\\
\text{for every integer }e\geq e_0\text{ and every }
 b\in\mathfrak a^{\lceil tp^e\rceil},\\
 cb(1\otimes z)=0\text{ in }F_R^e(M)
\end{array}\right\}.
\]
Here \(\mathfrak a^t\) denotes the ideal--exponent pair used in tight closure, rather than an ordinary ideal power with a nonintegral exponent. For a subset \(N\subseteq M\), write \(\operatorname{Ann}_R(N)=\{r\in R:rn=0\text{ for every }n\in N\}\). Under the stated convention, define
\[
\tau_{\mathrm{fg}}(\mathfrak a^t)
=\bigcap_{\substack{M\text{ a finitely generated}\\R\text{-module}}}
\operatorname{Ann}_R\bigl(0_M^{*\mathfrak a^t}\bigr),
\qquad
\tau_b(\mathfrak a^t)
=\bigcap_{M\text{ an }R\text{-module}}
\operatorname{Ann}_R\bigl(0_M^{*\mathfrak a^t}\bigr).
\]
The question is whether, for every \(F\)-finite Noetherian domain \(R\) of positive characteristic, every nonzero ideal \(\mathfrak a\subseteq R\), and every \(t\in\R_{>0}\), one has
\[
\tau_{\mathrm{fg}}(\mathfrak a^t)=\tau_b(\mathfrak a^t).
\]
No locality, normality, or further hypothesis is stated in the supplied record. The real-exponent interpretation and the explicit tight-closure conventions above require comparison with the cited source before acceptance.
\subsection{Short English statement}
The question asks whether testing the tight closure of zero in finitely generated modules gives the same ideal as testing it in all modules, for every \(F\)-finite Noetherian domain of positive characteristic, every nonzero ideal, and every positive exponent. This draft uses positive real exponents and explicitly stated standard conventions, which the supplied record does not itself define.
\subsection{Sources}
\begin{itemize}
\item Local UnsolvedMath registry, \texttt{data/unsolvedmath/problems.json}, permanent ID \(30004079\), title \emph{Equality of Finitistic and Big Test Ideals}; supplied source problem code \texttt{catalogue entry}. The record belongs to the Oberwolfach Reports Open Problems collection and gives proposed year 2019; no individual proposer is supplied.
\item Supplied source citation: \emph{Singularities and Homological Aspects of Commutative Algebra} (2020), Oberwolfach Reports, pp.~407--464, DOI \url{https://doi.org/10.4171/owr/2019/6}. The supplied background describes Oberwolfach Reports as records of workshops and their open-problem discussions.
\item Supplied literature reference: Oberwolfach Report 6/2019, \url{https://doi.org/10.4171/owr/2019/6}.
\item Supplied literature reference: Takagi, finitistic test ideals, \url{https://doi.org/10.1016/j.jalgebra.2018.08.003}.
\item Supplied literature reference: Aberbach--Huneke--Polstra, \url{https://doi.org/10.1016/j.aim.2024.109559}.
\item Literature assessment dated exactly \texttt{2026-08-22}: the supplied curation labels the problem open and reports equality in known cases under additional hypotheses, while describing the general conjecture as remaining open. This is supplied dated curation, not an independently verified assessment. The record requests independent verification of time-sensitive status and cited claims.
\item The supplied statement-status label is \texttt{corrected\_verified}; its accompanying note says that mathematical content was preserved while presentation or omitted definitions were repaired. This records the supplied editorial designation and does not establish mathematical verification of this draft.
\end{itemize}
\end{document}
